\section*{第 \ref{ch:b1:species-interactions} 章——物种间的相互作用}

\begin{solution}{exo:b1:species-interactions:1}
\omterm{def:b1:species-interactions:types}{竞争} $-/-$（两棵橡树争夺光照）；\omterm{def:b1:species-interactions:types}{捕食} $+/-$（猞猁、野兔）；
\omterm{def:b1:species-interactions:types}{寄生} $+/-$（蜱虫、鹿）；\omterm{def:b1:species-interactions:types}{互利共生} $+/+$（蜜蜂、花）；
\omterm{def:b1:species-interactions:types}{偏利共生} $+/0$（树皮上的蕨类）；偏害\omterm{def:b1:species-interactions:symbiosis}{共生} $-/0$（
橡树树荫下的草）。
\end{solution}

\begin{solution}{exo:b1:species-interactions:2}
两个以相同方式利用同一种有限资源的物种，不能
在稳定\omterm{def:b1:organism-environment:organism}{环境}中无限期共存。这五种莺都取食相同的食物，
但取食于树的不同部位，且取食时间也不同：
它们的\omterm{def:b1:ecosystem-organization:niche}{生态位}不同，因此它们之中没有任何两者以“相同方式”利用资源。
\end{solution}

\begin{solution}{exo:b1:species-interactions:3}
\omterm{def:b1:species-interactions:functional}{处理时间}：捕获、进食并消化一个猎物所需的时间，在此期间
不\omterm{def:b1:species-interactions:types}{捕食}其他猎物。攻击率：单位时间内搜索的面积
乘以捕获概率。最大速率 $1/h = 2$ 个猎物每
小时。
\end{solution}

\begin{solution}{exo:b1:species-interactions:4}
若移除某一物种会使\omterm{def:b1:ecosystem-organization:ecosystem}{群落}结构发生与其丰度不成比例的改变，
这种物种即为\omterm{def:b1:species-interactions:keystone}{关键种}，通常是通过控制某一占优势的
\omterm{def:b1:species-interactions:types}{竞争}者来实现的。丰度只衡量存在与否，而不衡量作用大小：一种稀少的\omterm{def:b1:species-interactions:types}{捕食}者
若\omterm{def:b1:species-interactions:types}{捕食}了本会胜出的物种，其作用便无法从数量上体现出来。
\end{solution}

\begin{solution}{exo:b1:species-interactions:5}
\omterm{prop:b1:species-interactions:lotka}{等斜线} 1：$N_1 + 0.5 N_2 = 500$（$N_1$ 轴上截距 500，
$N_2$ 轴上截距 1000）；\omterm{prop:b1:species-interactions:lotka}{等斜线} 2：$N_2 + 0.6 N_1 = 400$（$N_2$ 轴上截距 400，$N_1$ 轴上截距 667）。
两线相交，且每个物种受自身的限制都大于受对方的限制（$\alpha = 0.5 <
K_1/K_2 = 1.25$；$\beta = 0.6 < K_2/K_1 = 0.8$），两者稳定共存，
平衡点为 $N_1 = 429$，$N_2 = 143$。当 $\alpha = 2$ 时，\omterm{prop:b1:species-interactions:lotka}{等斜线} 1 的截距为
500 和 250，二者都位于\omterm{prop:b1:species-interactions:lotka}{等斜线} 2 的截距（667 和 400）之内：物种 2
总是获胜。
\end{solution}

\begin{solution}{exo:b1:species-interactions:6}
$f = aN/(1 + ahN)$：密度为 10 时，$2/(1 + 0.5) = 1.33$ 个每小时；密度为 100 时，
$20/(1 + 5) = 3.33$ 个每小时（最大值为 $1/h = 4$）。半饱和
密度为 $N = 1/(ah) = 20$ 个每平方米。
\end{solution}

\begin{solution}{exo:b1:species-interactions:7}
对于 III 型反应，低密度时被吃掉猎物的比例随密度
上升而增加：稀少的猎物被吃掉的比例相对较低，
因而能够恢复。对于 II 型反应，被吃掉的比例在低密度时
最高（此时\omterm{def:b1:species-interactions:types}{捕食}者从不饱和），因此小规模的猎物\omterm{def:b1:populations:population}{种群}
会被进一步压低。
\end{solution}

\begin{solution}{exo:b1:species-interactions:8}
\emph{P.~caudatum} 在同一水层中以与
\emph{P.~aurelia} 相同的细菌为食，而生长更快、耗食更省的 \emph{P.~aurelia}
把食物压到其所需水平以下。\emph{P.~bursaria} 则在底层
以不同的资源为食：决定结果的是\omterm{def:b1:ecosystem-organization:niche}{生态位}，而非生长速率
——\emph{P.~bursaria} 的生长速度并不比 \emph{P.~caudatum} 快。
\end{solution}

\begin{solution}{exo:b1:species-interactions:9}
\omterm{def:b1:species-interactions:types}{捕食}者更频繁地遇到这种斑纹而不受伤害，学习记住它的程度降低，
且更多地攻击它：拟态者与模式者的保护作用都随之下降。
因此，拟态者随数量增多而适合度下降，形成负频率依赖，
使其相对于模式者保持稀少。
\end{solution}

\begin{solution}{exo:b1:species-interactions:10}
没有\omterm{def:b1:species-interactions:types}{捕食}者时，贻贝的\omterm{prop:b1:species-interactions:lotka}{等斜线}位于每一个\omterm{def:b1:species-interactions:types}{竞争}者的\omterm{prop:b1:species-interactions:lotka}{等斜线}之外：
它将赢得这块岩石。\omterm{def:b1:species-interactions:types}{捕食}降低了贻贝的有效 $K$，
（其\omterm{prop:b1:species-interactions:lotka}{等斜线}向内收缩），直到从内侧穿过其他物种的\omterm{prop:b1:species-interactions:lotka}{等斜线}：
\omterm{def:b1:species-interactions:types}{竞争}者现在受自身的限制大于受贻贝的限制，
从而与贻贝共存。\omterm{def:b1:species-interactions:types}{捕食}者
作用于最强的\omterm{def:b1:species-interactions:types}{竞争}者，为其余物种腾出空间。
\end{solution}

\begin{solution}{exo:b1:species-interactions:11}
当磷酸盐限制生长时，\qty{80}{\%} 的磷酸盐
远比 \qty{15}{\%} 的糖更有价值：\omterm{def:b1:species-interactions:types}{互利共生}。在磷酸盐丰富的
土壤中，植物自己就能吸收磷酸盐，而真菌
白白索取 \qty{15}{\%}：\omterm{def:b1:species-interactions:types}{寄生}。实验：在一系列磷酸盐水平下，分别培养
有真菌与无真菌的植物，并测量
生物量；\omterm{def:b1:plant-water-minerals:mycorrhiza}{菌根}植物长得更大处为\omterm{def:b1:species-interactions:types}{互利共生}，
长得更小处为\omterm{def:b1:species-interactions:types}{寄生}。
\end{solution}

\begin{solution}{exo:b1:species-interactions:12}
围栏样地：沿河设置的围栏地块，麋鹿无法进入，
与开放样地配对，在若干山谷中重复设置，有狼群与无狼群两类，
在狼群重新引入前后连续测量多年；
记录柳树的高度与盖度、麋鹿的密度以及其在各样地中停留的时间
（通过足迹与照相机），以及狩猎与积雪深度。级联
假说预测：柳树仅在有狼的开放样地中恢复，
且恢复速度与围栏内相同；若柳树无论有无狼都同样恢复，
或其恢复随积雪与狩猎而变化，
则级联假说被否定，或与这些因素共同起作用。
\end{solution}

\begin{solution}{pb:b1:species-interactions:1}
\textbf{1.} A：$N_A + 1.4 N_B = 200$，截距为 200（$N_A$）与 143
（$N_B$）。B：$N_B + 0.9 N_A = 130$，截距为 130（$N_B$）与 144
（$N_A$）。
\textbf{2.} 截距（$N_A$ 轴上为 200 对 144，$N_B$ 轴上为 143
对 130）表明 A 的\omterm{prop:b1:species-interactions:lotka}{等斜线}完全位于 B 的\omterm{prop:b1:species-interactions:lotka}{等斜线}之外：
解 $N_A = 200 - 1.4 N_B$ 与 $N_B = 130 - 0.9 N_A$ 得 $N_A = -69$，
因此两直线的交点位于正象限之外。无论从何种初始状态出发，
只要 B 仍能增长，A 也增长，并且比 B 存活得更久。
\textbf{3.} $T_A = 0.87$ 天，$T_B = 1.26$ 天：A，
即增长更快的一方。
\textbf{4.} A：$N_A + 0.3 N_C = 200$（200、667）；C：$N_C + 0.3 N_A =
100$（100、333）。每个物种受自身的限制都大于受对方的限制
（$0.3 < 200/100$，且 $0.3 < 100/200 = 0.5$）：稳定共存。
\textbf{5.} 森林并非恒定不变（季节更替与干扰会在\omterm{def:b1:species-interactions:types}{竞争}分出胜负之前重新开始这一\omterm{def:b1:species-interactions:types}{竞争}），
并拥有众多资源与\omterm{def:b1:ecosystem-organization:ecosystem}{生境}，因此
物种会划分\omterm{def:b1:ecosystem-organization:niche}{生态位}，而不是共享同一个。
\textbf{6.} $N_A = 200 - 0.3 N_C$，$N_C = 100 - 0.3 N_A = 100 - 60 +
0.09 N_C$，故 $N_C = 44$，$N_A = 187$。
\textbf{7.} 上升并在约 10 处趋于平缓：II 型。
\textbf{8.} $(1/N, 1/f)$：$(0.2, 0.5)$，$(0.1, 0.303)$，$(0.05, 0.2)$，
$(0.025, 0.149)$，$(0.0125, 0.125)$：一条直线。
\textbf{9.} 斜率 $(0.5 - 0.125)/(0.2 - 0.0125) = 2.0 = 1/a$，故 $a =
\qty{0.5}{m^2/d}$；截距 $0.5 - 2\times 0.2 = 0.1$ 天 $= h$。
\textbf{10.} 最大值 $1/h = 10$ 个猎物每天；半饱和密度 $N =
1/(ah) = 20$ 个每平方米。
\textbf{11.} 处理比例 $= f h = f/10$：当 $f
> 5$ 时超过一半，即在每平方米 40 和 80 个时（在 20 时恰为 5.0，即一半）。
\textbf{12.} 猎物增长 $0.1 N (1 - N/200)$ 等于\omterm{def:b1:species-interactions:types}{捕食}量
$0.5N/(1 + 0.05N)$：$0.1(1 - N/200) = 0.5/(1 + 0.05 N)$，即
$(1 - N/200)(1 + 0.05N) = 5$：$1 + 0.05N - N/200 - N^2/4000 = 5$，故
$N^2 - 180 N + 16\,000 = 0$，$N = 90 \pm \sqrt{8100 - 16\,000}$：无
实根——每平方米一只甲虫的\omterm{def:b1:species-interactions:types}{捕食}量超过了猎物所能产生的数量
（猎物的最大增长 $rK/4 = 5$ 每天，出现在 $N = 100$ 处，此时
甲虫吃掉 8.3 个），并将其导向灭绝。不存在平衡点；
猎物只有在甲虫较少或拥有避难所时才能存续。
\textbf{13.} 对第二种猎物，\omterm{def:b1:species-interactions:functional}{处理时间}减半（$h' = 0.05$），
因此在饱和时，甲虫每天最多可吃掉其中 20 个：总摄入量
在高密度下若转而\omterm{def:b1:species-interactions:types}{捕食}处理更快的猎物，大致翻一番。
\textbf{14.} $H_{15} = \ln 15 = 2.71$。八个物种：$-(0.85\ln 0.85
+ 7\times 0.0214\ln 0.0214) = 0.138 + 0.576 = 0.71$。
\textbf{15.} 相差的倍数为 $2.71/0.71 = 3.8$。
\textbf{16.} 海星对贻贝的 $-$ 作用，使贻贝对其\omterm{def:b1:species-interactions:types}{竞争}者的 $-$ 作用
减弱；移除海星后，贻贝便达到其\omterm{def:b1:species-interactions:types}{竞争}能排斥其余物种的
密度。
\textbf{17.} 生物量衡量的是某一物种占据了多少能量流，
而不是其相互作用的效应；\omterm{def:b1:species-interactions:keystone}{关键种}的作用是通过它所控制的
物种实现的，而该物种的生物量很大。
\textbf{18.} 若移除海星，将释放一个稀少的物种，
该物种无法占据这块岩石：以贻贝为主的\omterm{def:b1:ecosystem-organization:ecosystem}{群落}
几乎不会改变，海星也就不会是
\omterm{def:b1:species-interactions:keystone}{关键种}。
\textbf{19.} 用笼子将海星排除在小片样地之外，以开放的笼子
作为对照，以排除笼子本身的作用：若贻贝只在封闭的笼子中
占据优势，则对贻贝的\omterm{def:b1:species-interactions:types}{捕食}即为其机制；
若将海星留在笼中而人工移除贻贝，
则可检验海星的其他猎物是否重要。
\textbf{20.} 获得 $\qty{40}{mg}$，代价为 $8\times 0.5 = \qty{4}{mg}$：净得
\qty{36}{mg} $=$ \qty{0.58}{kJ}。
\textbf{21.} \qty{40}{\micro g} 对应 \qty{4}{mg}：花每日光合产物的 \qty{1}{\%}
可换取约三次花粉传递，这即
植物的繁殖。
\textbf{22.} 双方都从对方取得某些东西，但各自所得多于所失，
符号取决于净效应，而不取决于
利用本身。
\textbf{23.} 盗蜜者为 $+$，植物为 $-$（花蜜被盗而未传粉）：这是对\omterm{def:b1:species-interactions:types}{互利共生}的一种
\omterm{def:b1:species-interactions:types}{寄生}；蜜蜂发现花已被采空，所获更少，
是间接的 $-$。
\textbf{24.} 植物 $\leftrightarrow$ 蜜蜂 $+/+$；盗蜜者 $\to$ 植物 $-$，
盗蜜者 $\to$ 蜜蜂 $-$；鸟 $\to$ 盗蜜者 $-$。两个负号：鸟
对植物和蜜蜂都有间接的 $+$ 作用。
\textbf{25.} $a = \qty{0.5}{m^2/d}$，$h = \qty{0.1}{d}$（\qty{2.4}{h}），
最大值 \num{10} 个猎物每天。
\end{solution}
